\begin{table*}[p]
    \centering
    \resizebox{\textwidth}{!}{
    \setlength{\tabcolsep}{2pt}
    \begin{tabular}{lccccccccc}
    \toprule
          \rotatebox{45}{Analysis} 
        & \rotatebox{45}{Evolutionary VCV}
        & \rotatebox{45}{Modern human VCV}
        & \rotatebox{45}{Neandertal VCV}
        & \rotatebox{45}{\textit{Pan paniscus} VCV}
        & \rotatebox{45}{\textit{Pan troglodytes} VCV}
        & \rotatebox{45}{\textit{G. beringei} VCV}
        & \rotatebox{45}{\textit{G. gorilla} VCV}
        & \rotatebox{45}{\textit{Po. abelii} VCV}
        & \rotatebox{45}{\textit{Po. pygmaeus} VCV} \\
    \midrule
    C\textsubscript{1} & 23.68 (13.21) & 108.06 (27.99) & 43.91 (16.33) & 94.79 (14.4) & 61.95 (16.25) & 29.91 (11.34) & 13.31 (8.86) & 22.91 (11.55) & 31.74 (8.73) \\
    \midrule
    I\textsuperscript{2} & 12.38 (9.42) & 121.87 (47.96) & 40.85 (14.45) & 12.02 (10.18) & 46.47 (13.44) & - & - & - & - \\
    \bottomrule
    \end{tabular}}
    \caption{Kullback-Leibler divergences from the prior to the posterior distribution for the variance-covariance matrices. These posterior distributions are Inverse Wishart distributed; the maximum likelihood estimate for each posterior distribution's degree of freedom is listed in parentheses. Lower degrees of freedom indicate heavier tails.}
    \label{table:klDivergences}
\end{table*}
